\chapter[Global Diffeomorphism]{Global Diffeomorphism}
Throughout this chapter let
\[
F = Q \circ L :
\Delta_k^\circ
\longrightarrow
\operatorname{Im}(Q)
\]
denote the canonical logarithmic coordinate map introduced in Chapter 4.
Chapter 5 established that (F) is a local diffeomorphism. The objective
of the present chapter is to prove that (F) is globally injective and
consequently a global diffeomorphism onto its image.
\section{Global Injectivity}
\begin{theorem}[Global Injectivity]
\label{ch:thm:6-1}
The canonical logarithmic coordinate map
\[
F :
\Delta_k^\circ
\longrightarrow
\operatorname{Im}(Q)
\]
is injective.
\end{theorem}
\begin{proof}
Assume
\[
\mathcal P
=
\prod_{c\in\mathcal C}
\Delta_c^\circ.
\]
By definition,
\[
F((P_c)_c)
=
(Q(\log P_c))_{c\in\mathcal C}.
\]
Since
\[
F :
\mathcal P
\longrightarrow
\prod_{c\in\mathcal C}\operatorname{Im}(Q_c)
\]
there exists a scalar
\[
F :
\mathcal P
\longrightarrow
F(\mathcal P)
\]
such that
Hence, for every coordinate,
Exponentiating,
Summing over all coordinates,
Therefore
Since the exponential function is injective,
Consequently,
and therefore
Hence (F) is injective. 
\end{proof}
\section{Product Posterior Manifold}
The posterior manifold associated with a Weighted Incidence Structure is
\[
\mathcal P
=
\prod_{c\in\mathcal C}
\Delta_c^\circ.
\]
The canonical logarithmic coordinate map acts componentwise,
\[
F((P_c)_c)
=
(Q(\log P_c))_{c\in\mathcal C}.
\]
\begin{corollary}[Product Injectivity]
\label{ch:cor:6-2}
The canonical logarithmic coordinate map
\[
F :
\mathcal P
\longrightarrow
\prod_{c\in\mathcal C}\operatorname{Im}(Q_c)
\]
is injective.
\end{corollary}
\begin{proof}
Suppose
Then
for every ciphertext block (c).
Applying Theorem~\ref{ch:thm:6-1} to each block yields
Therefore
Hence (F) is injective. 
\end{proof}
\section{Global Diffeomorphism}
\begin{theorem}[Global Diffeomorphism]
\label{ch:thm:6-3}
The canonical logarithmic coordinate map
\[
F :
\mathcal P
\longrightarrow
F(\mathcal P)
\]
is a diffeomorphism onto its image.
\end{theorem}
\begin{proof}
By Proposition~\ref{ch:prop:5-1}, the map (F) is smooth.
By Theorem~\ref{ch:thm:5-4}, (F) is a local diffeomorphism.
By Corollary~\ref{ch:cor:6-2}, (F) is injective.
A standard theorem in differential topology (see, e.g., Lee,
*Introduction to Smooth Manifolds*, theorem on injective local
diffeomorphisms) states that an injective local diffeomorphism is a
diffeomorphism onto its image.
Therefore
is a global diffeomorphism. 
\end{proof}
\section{Consequences}
Theorem~\ref{ch:thm:6-3} provides a globally defined logarithmic coordinate
representation of the posterior manifold modulo additive gauge
transformations.
Consequently, every differential-geometric construction introduced in
the subsequent chapters may be performed either on the posterior
manifold or on its logarithmic image without loss of information.
The pullback metric, curvature, and the Universal Optimality Defect will
therefore be developed on equivalent geometric representations of the
same posterior structure.
